\section{Valorisation scientifique des collections}
L'éditorialisation constitue désormais un maillon indispensable pour donner de la cohérence intellectuelle à une bibliothèque numérique. L'éditorialisation contextualise les fonds, guide l'usager et valorise l'exxpertise scientifique de l'établissement.

Lors des enquêtes de terrains, les chercheurs ont fortement insisté sur la mise en place de cette médiation. L'un des chercheurs souligne qu'actuellement dans CujasNum, cela manque parfois de textes explicatifs, notamment sur la constitution des collections ou explicitant les choix de classifications. De même, un autre professeur rappelle l'importance d'attirer l'attention des usagers sur le contenu et la valeur singulière des ouvrages anciens.

Pour mettre en place cette médiation culturelle sans alourdir l'infrastructure technique, la bibliothèque Cujas peut s'inspirer du fonctionnement de la BIS. Lors de la conception d'expositions virtuelles ou de billets d'actualité, la BIS renonce à l'utilisation d'exposition virtuelles intégrées dans des visionneuses lourdes et peu portables. Elle a donc décidé d'opter pour le format Markdown. Les billets de blog et les pages d'expositions deviennent de simples fichiers textes structurés par des balises Markdown très légères, et exportables d'une solution technique à une autre. Ce choix technique permet de s'affranchir des instabilités des extensions de CMS, assurant une pérennité des données.

Pour rationaliser et homogénéiser la production éditoriale scientifique écrite par des tiers, la plateforme doit appliquer une démarche de qualité rigoureuse inspirée du guide de NuBIS. ce guide impose des règles strictes autour de trois points : la gestion des images, les contraintes rédactionnelles et le respect des droits d'auteur. Le respect de ces trois règles imposent un cadre stricte que doivent suivre les chercheurs souhaitant participer à l'éditorialisation de leur bibliothèque numérique. Comme l'a expliqué Elina Leblanc qui a analysé les \enquote{expositions nativement numériques} comme des espaces collaboratifs de recherche. Elle parle des avantages de permettre à des chercheurs, des enseignants ou des doctorants de rédiger des focus ou des notices scientifiques \enquote{participe au processus d'éditorialisation et conduit ainsi à créer de nouveaux savoirs, pouvant à leur tout être partagés}.

\subsection{Les formats de valorisation}
En s'appuyant sur ces règles de rédaction, l'éditorialisation scientifique de la BIU Cujas se déploiera sous plusieurs formes complémentaires, penser pour différents niveaux de lecture, comme la Galerie des grands juristes et l'éditorialisation par des juristes pour des juristes.

La \enquote{Galerie des grands juristes} est proposé par le directeur de la bibliothèque Cujas. Cette entrée thématique prendra la forme d'une galerie de portraits intéractifs. Inspirée de la même idéée que la bibliothèque de l'ENS qui a mis en avant ces collections, sous le nom des \textit{Parcours normalien.ne.s}. Chaque juriste fera l'objet d'une notice biographique, compris entre 5 000 et 10 000 signes, rédigée par un chercheur. Depuis la fiche bibliographique de l'auteur, l'usager pourra rebondir instantanément sur l'ensemble de ses oeuvres numérisées disponible dans la bibliothèque Cujas ou sur des sites partenaires.

Une deuxième forme d'éditorialisation des collections sera d'écrire des billets de blog, des parcours thématiques, écrite par des spécialistes d'une question ou d'une thématique pour d'autre s'intéressant au sujet. Cela pourrait également comme le suggérait l'un des enseignats-chercheurs une porte d'entrée pour les jeunes chercheurs pour développer leur carrière. L'ancienne responsable de la bibliothèque numérique rappelle l'intérêt d'associer directement les étudiants avancés et les chercheurs à cette démarche de valorisation. L'idée serait de permettre à des chercheurs ou des étudiants d'écrire des articles des billets de blog sur un ou des documents étudiés pour leur étude provenant de la bibliothèque Cujas à le replacer dans un contexte intellectuel, redonnant \enquote{vie} aux documents numérisés. Laurent Pfister pense que ce serait un excellent moyen de valoriser la production scientifique notamment pour des jeunes chercheurs.

En combinant la flexibilité technique du format Markdown et la rigueur des guides d'expositions de la BIS, la bibliothèque numérique de Cujas ne sera plus uniquement un outil d'exposition de ses documents numérisés, mais deviendrait une véritable plateforme de médiation culturelle juridique. POur que ces services soient utilisés par les usagers, il est nécessaire pour les experts métiers de les accompagner dans leur prise en main. Il est important alors que les experts métiers eux-mêmes soient capables de les utiliser amenant le service dédié de la bibliothèque à les accomapgner dans la conduite du changement.